%%
%% spring-lecture11.tex
%% 
%% Made by Alex Nelson <pqnelson@gmail.com>
%% Login   <alex@lisp>
%% 
%% Started on  2026-04-23T12:49:06-0700
%% Last update 2026-04-23T12:49:06-0700
%% 

\lecture[Variational properties of geodesics]

\begin{theorem}
Let $(M,g)$ be a Riemannian manifold. Then the following facts hold:
\begin{enumerate}
\item $\exp\colon\Omega\subset TM\to M$ is smooth
\item For every $x\in M$, $\D(\exp_{x})_{0}=\id_{T_{x}M}$
\item The map $\Phi\colon\Omega\to M\times M$ given by
  $\Phi(x,v)=(x,\exp_{x}(v))$ is a local diffeomorphism at $(x,0)$ for
  every $x\in M$.
\end{enumerate}
\end{theorem}

\begin{proof}
\begin{enumerate}
\item This follows immediately from the local existence and uniqueness
  of geodesics.
\item If this is the identity, then $\exp_{x}\colon T_{x}M\to M$, but
\begin{equation}
\D(\exp_{x})_{0}\colon T_{0}(T_{x}M)\to T_{\exp_{x}(0)}M=T_{x}M
\end{equation}
since $\exp_{x}(0)=x$ and $T_{0}(T_{x}M)\iso T_{x}M$ canonically, so
$\D(\exp_{x})_{0}$ can be viewed as an endomorphism of $T_{x}M$.
We really have $\exp_{x}\colon V_{x}\to M$ where $V_{x}\subset T_{x}M$
is an open set.

Choose any $v\in V_{x}$ and let $s\in(0,1)$ be sufficiently small such
that $sv\in V_{x}$. Let $I\subset\RR$ be the domain of $\gamma_{x,v}$
(the maximal geodesic with $\gamma_{x,v}(0)=x$ and $\gamma'_{x,v}(0)=v$).
Define the interval $\widetilde{I}=\{t\in\RR\mid st\in I\}$. Note that
$1\in\widetilde{I}$ (we know $1\in I$ since $st\in I$ so $1\in\widetilde{I}$),
and that the smooth curve
\begin{equation}
\alpha_{x,v,s}\colon\widetilde{I}\to M
\end{equation}
given by $\alpha_{x,v,s}(t)=\gamma_{x,v}(st)$ for every
$st\in\widetilde{I}$ is a geodesic with initial conditions
$\alpha_{x,v,s}(0)=x$ and
$\alpha'_{x,v,s}(0)=sv=s\gamma'_{x,v}(0)$. Therefore,
\begin{equation}
\alpha_{x,v,s}(t)=\gamma_{x,sv}(t)
\end{equation}
for every $t\in\widetilde{I}$. (Then we found how
$\gamma_{x,v}(st)=\gamma_{x,sv}(t)$ rescaling from $I$ to
$\widetilde{I}$ rescales the initial condition for velocity from $v$
to $sv$.)

Now, we're ready! Then we have
\begin{subequations}
  \begin{align}
\D(\exp_{x})_{0}(v) &= \left.\frac{\D}{\D s}\exp_{x}(sv)\right|_{s=0}\\
&=\left.\frac{\D}{\D s}\gamma_{x,sv}(1)\right|_{s=0}\\
&=\left.\frac{\D}{\D s}\alpha_{x,v,s}(1)\right|_{s=0}\\
&=\left.\frac{\D}{\D s}\gamma_{x,v}(s)\right|_{s=0}\\
&=\gamma'_{x,v}(0)=v.
  \end{align}
\end{subequations}
Hence $\D(\exp_{x})_{0}$ is the identity.
\item We have $\Phi\colon\Omega\subset TM\to M\times M$ and
\begin{equation}
\D\Phi_{(x,0)}\colon T_{(x,0)}\Omega\to T_{(x,x)}(M\times M)\iso T_{x}M\times T_{x}M.
\end{equation}
So $\D\Phi_{(x,0)}$ is an endomorphism of $T_{x}M\times T_{x}M$. For
every $(w_{1},w_{2})\in T_{x}M\times T_{x}M$, consider a curve
\begin{equation}
(x,v)\colon(-\varepsilon,\varepsilon)\to\Omega
\end{equation}
such that
\begin{subequations}
\begin{equation}
(x,v)(0)=(x(0),v(0))=(x,0)
\end{equation}
and
\begin{equation}
(x,v)'(0)=(x'(0),v'(0))=(w_{1},w_{2}).
\end{equation}
\end{subequations}
Then
\begin{subequations}
  \begin{align}
\D\Phi_{(x,0)}(w_{1},w_{2})
&=\left.\frac{\D}{\D t}\Phi(x(t),v(t))\right|_{t=0}\\
&=\left.\frac{\D}{\D t}\bigl(x(t),\exp_{x(t)}(v(t))\bigr)\right|_{t=0}\quad\mbox{by def.~of }\Phi\\
&=\left(\left.\frac{\D x(t)}{\D t}\right|_{t=0},\left.\frac{\D}{\D t}\exp_{x(t)}(v(t))\right|_{t=0}\right)\\
&=\left(x'(0),\left.\frac{\D}{\D t}\exp(x(t),v(t))\right|_{t=0}\right)\\
&=\left(x'(0),\left.\frac{\D}{\D t}\exp(x(t),v(0))\right|_{t=0}+\left.\frac{\D}{\D t}\exp(x(0),v(t))\right|_{t=0}\right)\\
&=(x'(0),x'(0)+\D(\exp_{x})_{0}(w_{2}))\\
&=(w_{1},w_{1}+w_{2}).
  \end{align}
\end{subequations}
Since this is invertible, it follows that $\Phi$ is a local diffeomorphism.
  \qedhere
\end{enumerate}
\end{proof}

The third item from the previous theorem has an immediately
consequence:

\begin{corollary}
Let $(M,g)$ be a Riemannian manifold. For every $x_{0}\in M$, there
exists $\varepsilon>0$ and an open neighborhood $U\subset M$ of
$x_{0}\in U$ such that the following all hold:
\begin{enumerate}
\item For every $x,y\in U$ there exists a unique $v_{y}\in T_{x}M$
  such that $|v_{y}|_{g}<\varepsilon$ and $\exp_{x}(v_{y})=y$
\item The map $(x,y,t)\mapsto\gamma_{x,v_{y}}(t)$ is smooth where defined
\item For every $x\in U$, the map $\exp_{x}\colon B_{\varepsilon}(0)\subset T_{x}M\to M$
  is a diffeomorphism onto its image.
\end{enumerate}
\end{corollary}

\begin{proof}
Exercise.
\end{proof}

\begin{node}
It would be nice to speak of ``open balls'' in a manifold in general,
since open balls are the nicest guys in topology. We have a notion of
distance to construct balls on manifolds, but that's \emph{not} the
``good way'' to construct a \emph{useful} notion of an open ball. No!
The ``correct'' way is to look at (3) of this previous corollary and
use the image of an open ball $B_{\varepsilon}(0)\subset T_{x}M$ as
the open ball in $M$.
\end{node}

\begin{definition}
Let $(M,g)$ be a Riemannian manifold. Let $x\in M$.
\begin{enumerate}
\item The \define{Injectivity Radius} is the number
  $\inj_{x}(M)=\sup\{r>0\mid \exp_{x}\colon B_{r}(0)\to M\mbox{ is a diffeomorphism onto its image}\}$.
  This is a positive number.
\item For every $r\in(0,\inj_{x}(M))$, the \define{Geodesic Ball} at
  $x$ of radius $r$ is the open subset
  $B^{M}_{r}(x):=\exp_{x}(B_{r}(0))\subset M$.
\item For every $r\in(0,\inj_{x}(M))$, the \define{Geodesic Sphere} at
  $x$ of radius $r$ is the open subset
  $S^{M}_{r}(x):=\exp_{x}(\boundary B_{r}(0))=\boundary B^{M}_{r}(x)\subset M$.
\end{enumerate}
\end{definition}

\subsection{Variational Properties of Geodesics}

\begin{node}
What are ``variational properties''? Well, geodesics arise as critical
points of functionals of curves on manifolds. Length is not a ``nice''
functional due to the presence of the squareroot in the integrand. We
construct a related (but nicer) functional.
\end{node}

\begin{definition}
Let $(M,g)$ be a Riemannian manifold. Let $\gamma\colon[a,b]\to M$
be a smooth curve. The \define{Energy} of $\gamma$ is given by the
functional
\begin{equation}
E_{g}(\gamma):=\frac{1}{2}\int^{b}_{a}g_{\gamma(t)}(\gamma'(t),\gamma'(t))\,\D t.
\end{equation}
\end{definition}

\begin{node}
Now, what does it mean for $\gamma$ to be a ``critical point'' of
$E_{g}$? Well, its first variation must vanish.
\end{node}

\begin{definition}
Let $(M,g)$ be a Riemannian manifold. Let $\gamma\colon[a,b]\to M$ be a
smooth curve. We say $V\in\Vect{\gamma}$ is a \define{Variational Field}
if there exists a smooth map
$F\colon[a,b]\times(-\varepsilon,\varepsilon)\to M$
such that
\begin{enumerate}
\item For any $s\in(-\varepsilon,\varepsilon)$, the map
  $\gamma_{s}:=F(-,s)\colon[a,b]\to M$ is a smooth curve on $M$ (so
  $\gamma_{s}(t)=F(t,s)$) with $\gamma_{s}(a)=\gamma(a)$ and
  $\gamma_{s}(b)=\gamma(b)$ the endpoints fixed
\item $\gamma_{0}=\gamma$ --- so $\gamma_{s}$ is a variation of $\gamma$
\item We have for every $t\in[a,b]$,
  \begin{equation}
\begin{split}
V(t) &:= \left.\frac{\partial}{\partial s}F(t,s)\right|_{s=0}\\
&=\D F_{(t,0)}\bigl((0,1)\bigr).
\end{split}
  \end{equation}
\end{enumerate}
\end{definition}